% ------------------------------------------------------------------------------
\section{Diseño de concreto y laboratorio}
\label{sec:con}
% ------------------------------------------------------------------------------

Este capítulo documenta el módulo de concreto: la biblioteca de materiales, la granulometría con husos
NTP~400.037 / ASTM~C33, el balance de volúmenes absolutos, el diseño de mezclas por el método ACI~211.1 con los
requisitos de la NTE~E.060, la corrección por humedad de una tanda, el registro de probetas con su evaluación y el
reporte en PDF. Las fórmulas y las tablas son las que contiene el código; donde el valor depende de una edición
concreta de una norma, se indica que debe verificarse.

% ------------------------------------------------------------------------------
\subsection{Flujo de datos entre pantallas}
\label{sec:con-flujo}

Las cuatro pantallas del módulo comparten la biblioteca de materiales y el proyecto activo
(\ruta{SesionProyecto}). Las mezclas y sus ensayos pertenecen a un proyecto; los materiales y sus granulometrías no.

\begin{table}[H]
	\centering
	\caption{Pantallas, controladores y servicios del módulo de concreto.}
	\label{tab:con-pantallas}
	\begin{tabular}{L{3.2cm}L{3.8cm}L{6.6cm}}
		\hline
		\textbf{Pantalla} & \textbf{Controlador} & \textbf{Servicios y cálculo} \\ \hline
		Materiales y ensayos & \ruta{MaterialesController} & \ruta{MaterialService} \\
		Granulometría & \ruta{GranulometriaController} & \ruta{GranulometriaService}, \ruta{HusosGranulometricos},
		\ruta{RegistroGranulometriaService} \\
		Balance de volúmenes & \ruta{BalanceController} & \ruta{VolumenAbsolutoService} \\
		Mezclas y probetas & \ruta{MezclasController} & \ruta{DisenoMezclaACI}, \ruta{MezclaService},
		\ruta{EvaluacionResistencia}, \ruta{ReporteMezcla} \\
		\hline
	\end{tabular}
\end{table}

El recorrido habitual es: registrar los agregados y el cemento con sus ensayos; verificar la granulometría de cada
agregado; obtener una dosificación por m$^3$ con el asistente ACI o con el balance de volúmenes; pasarla con «Usar en
una mezcla» al formulario «Registrar dosificación»; calcular la tanda corregida por humedad; registrar las probetas y
revisar la evaluación E.060; y guardar el PDF.

% ------------------------------------------------------------------------------
\subsection{Biblioteca de materiales}
\label{sec:con-materiales}

\ruta{service/MaterialService} guarda los materiales (entidad \ruta{Material}). Los tipos admitidos están en
\ruta{YEFRI\_LERMA\_TIPOS}, que también llena el combo de la pantalla:

\begin{table}[H]
	\centering
	\caption{Tipos de material (\ruta{MaterialService.YEFRI\_LERMA\_TIPOS}).}
	\label{tab:con-tipos}
	\begin{tabular}{L{4.4cm}L{9.2cm}}
		\hline
		\textbf{Tipo} & \textbf{Referencia indicada en el código} \\ \hline
		\ruta{CEMENTO} & Genérico de versiones anteriores, sin tipo. \\
		\ruta{CEMENTO TIPO I}, \ruta{II}, \ruta{V} & Portland, NTP~334.009. \\
		\ruta{CEMENTO TIPO IP}, \ruta{ICo} & Adicionados, NTP~334.090. \\
		\ruta{CEMENTO TIPO MS}, \ruta{HS} & Por desempeño, NTP~334.082. \\
		\ruta{AGREGADO FINO}, \ruta{AGREGADO GRUESO} & Agregados. \\
		\ruta{ADITIVO} & NTP~334.088 / ASTM~C494; la dosis se anota en observaciones. El diseño ACI de la aplicación no
		considera aditivos. \\
		\hline
	\end{tabular}
\end{table}

\begin{table}[H]
	\centering
	\caption{Validaciones de \ruta{MaterialService.\allowbreak{}yefriLermaGuardar}.}
	\label{tab:con-material-validacion}
	\begin{tabular}{L{5.4cm}L{8.2cm}}
		\hline
		\textbf{Campo} & \textbf{Regla} \\ \hline
		Nombre & Obligatorio, de 1 a 120 caracteres tras recortar espacios. \\
		Tipo & Uno de la tabla~\ref{tab:con-tipos}. \\
		Densidad relativa, módulo de finura, tamaño máximo (mm), peso unitario compactado (kg/m$^3$) & Opcionales; si se
		indican, finitos y mayores que cero. \\
		Absorción y humedad total (\%) & Opcionales; si se indican, entre 0 y 100. \\
		Fecha del ensayo & Opcional; si se indica, fecha ISO (AAAA-MM-DD). \\
		Procedencia, referencia del ensayo, observaciones & Hasta 255 caracteres cada uno. \\
		\hline
	\end{tabular}
\end{table}

\ruta{MaterialesController} lista los materiales ordenados por nombre, con un paginador, y muestra el detalle del
seleccionado. La pantalla \textbf{no edita ni elimina} materiales: cada resultado nuevo se registra como otro material,
para conservar la trazabilidad de los anteriores. Los mensajes de error del servicio se asocian al campo
correspondiente buscando palabras clave en el texto del mensaje (\ruta{yefriLermaCampoDelError}).

Otros servicios consultan la biblioteca con \ruta{yefriLermaPorTipo(tipo)} y \ruta{yefriLermaCementos()}, que devuelve
todos los tipos que empiezan por \ruta{CEMENTO}.

% ------------------------------------------------------------------------------
\subsection{Granulometría}
\label{sec:con-granulometria}

\subsubsection{Entrada de datos}

Las mediciones se escriben como texto, una línea por tamiz con el formato «abertura en mm; masa retenida en g», de
mayor a menor abertura. \ruta{RegistroGranulometriaService.\allowbreak{}yefriLermaLeer(String)} acepta coma decimal, ignora las
líneas vacías y rechaza textos vacíos o de más de 10\,000 caracteres y líneas que no tengan exactamente dos partes.
Aparte se ingresan la masa seca inicial $M_0$ y la masa en el fondo $F$.

\subsubsection{Cálculo}

\ruta{GranulometriaService.\allowbreak{}yefriLermaCalcular(M0, tamices, F)} produce una fila por tamiz. Con $R_i$ la masa retenida
en el tamiz $i$ (ordenados de mayor a menor abertura):

\begin{align}
	r_i &= \frac{R_i}{M_0} \cdot 100 && \text{(\% retenido)}, \\
	A_i &= \frac{1}{M_0} \sum_{k \le i} R_k \cdot 100 && \text{(\% retenido acumulado)}, \\
	P_i &= 100 - A_i && \text{(\% que pasa)}, \\
	M_r &= \sum_i R_i + F, \qquad \Delta = M_r - M_0, \qquad \Delta_{\%} = \frac{\Delta}{M_0} \cdot 100
	&& \text{(balance de masas)}.
\end{align}

Los porcentajes se refieren a la masa \textbf{inicial}, no a la recuperada. El cálculo rechaza $M_0 \le 0$, $F < 0$,
aberturas no positivas, repetidas o desordenadas, masas retenidas negativas y una masa acumulada mayor que $M_0$. La
diferencia $\Delta$ no se limita: se informa para que usted la compare con el procedimiento del laboratorio.

\subsubsection{Módulo de finura}

\ruta{HusosGranulometricos.\allowbreak{}yefriLermaModuloFinura} suma los porcentajes retenidos acumulados en la serie estándar

\begin{equation}
	\mathrm{MF} = \frac{1}{100} \sum_{a \in S} A(a), \qquad
	S = \{0.15;\ 0.30;\ 0.60;\ 1.18;\ 2.36;\ 4.75;\ 9.5;\ 19;\ 37.5;\ 75;\ 150\}\ \text{mm}.
\end{equation}

Un tamiz ingresado de abertura $a_i$ se identifica con la abertura $a$ de la serie si
$\lvert a_i - a \rvert \le 0{,}02\,a + 0{,}005$\,mm. Si se encuentran menos de seis tamices de la serie, el método
devuelve \ruta{null} (sin módulo de finura).

\subsubsection{Husos granulométricos}

\begin{table}[H]
	\centering
	\caption{Límites de \% que pasa de \ruta{HusosGranulometricos.Huso}.}
	\label{tab:con-husos}
	\begin{tabular}{cL{2.6cm}L{2.6cm}L{2.6cm}L{2.6cm}}
		\hline
		\textbf{Tamiz (mm)} & \textbf{Fino} & \textbf{Grueso N.º 57} & \textbf{Grueso N.º 67} & \textbf{Grueso N.º 7} \\ \hline
		37,5 & & 100 & & \\
		25 & & 95--100 & 100 & \\
		19 & & & 90--100 & 100 \\
		12,5 & & 25--60 & & 90--100 \\
		9,5 & 100 & & 20--55 & 40--70 \\
		4,75 & 95--100 & 0--10 & 0--10 & 0--15 \\
		2,36 & 80--100 & 0--5 & 0--5 & 0--5 \\
		1,18 & 50--85 & & & \\
		0,60 & 25--60 & & & \\
		0,30 & 5--30 & & & \\
		0,15 & 0--10 & & & \\
		\hline
	\end{tabular}
\end{table}

\ruta{yefriLermaVerificar(huso, filas)} busca cada tamiz del huso entre los ingresados con la misma tolerancia que el
módulo de finura y comprueba $P_{\min} \le P \le P_{\max}$ (con holgura de $10^{-9}$). El resultado:

\begin{itemize}
	\item \ruta{cumple} es verdadero solo si todos los tamices encontrados están dentro \textbf{y} no falta ninguno;
	\item si faltan tamices, el mensaje advierte que la verificación es parcial;
	\item si ningún tamiz coincide, el resultado es «no cumple» con el aviso correspondiente;
	\item para el huso fino incluye el módulo de finura y si está dentro del rango 2,3--3,1.
\end{itemize}

\subsubsection{Flujo de la pantalla}

\ruta{GranulometriaController}:

\begin{enumerate}
	\item El combo «Agregado de la biblioteca» lista solo materiales cuyo tipo empieza por \ruta{AGREGADO}. Al elegir
	      uno se carga su historial de ensayos guardados; elegir un ensayo del historial copia sus datos al formulario.
	\item «Calcular resultados» llena la tabla (abertura, retenido en g y \%, acumulado y pasa) y la curva granulométrica,
	      e informa la masa recuperada y la diferencia. Cualquier cambio en las masas o las mediciones borra el resultado
	      mostrado.
	\item La verificación depende del agregado elegido: si es \ruta{AGREGADO FINO}, se aplica el huso fino; si es
	      grueso, se prueban en orden los husos N.º 57, 67 y 7 y se informa el primero que se cumple por completo (o que
	      no cumple ninguno); sin agregado, solo se informa el módulo de finura.
	\item «Guardar ensayo» exige un agregado y guarda un \ruta{RegistroGranulometria} (material, identificación de
	      hasta 120 caracteres, fecha ISO, $M_0$, $F$ y el texto de mediciones) después de recalcularlo para validarlo.
\end{enumerate}

La curva es un gráfico de líneas con las aberturas como categorías, no en escala semilogarítmica. El módulo de finura
calculado no se copia al material: el valor que usa el diseño es el que se registró en la biblioteca.

\subsubsection{Ejemplo verificado}

La prueba \ruta{DisenoMezclaTest.\allowbreak{}husoDelAgregadoFinoYModuloDeFinura} usa una arena de 500\,g con 25\,g en el fondo:

\begin{table}[H]
	\centering
	\caption{Granulometría de la arena de prueba.}
	\label{tab:con-ejemplo-granulometria}
	\begin{tabular}{ccccc}
		\hline
		\textbf{Tamiz (mm)} & \textbf{Retenido (g)} & \textbf{Acumulado (\%)} & \textbf{Pasa (\%)} & \textbf{Huso fino} \\ \hline
		9,5 & 0 & 0 & 100 & 100 \\
		4,75 & 25 & 5 & 95 & 95--100 \\
		2,36 & 70 & 19 & 81 & 80--100 \\
		1,18 & 85 & 36 & 64 & 50--85 \\
		0,60 & 100 & 56 & 44 & 25--60 \\
		0,30 & 120 & 80 & 20 & 5--30 \\
		0,15 & 75 & 95 & 5 & 0--10 \\
		\hline
	\end{tabular}
\end{table}

\begin{equation}
	\mathrm{MF} = \frac{0 + 5 + 19 + 36 + 56 + 80 + 95}{100} = 2.91 .
\end{equation}

La arena cumple el huso fino y su módulo de finura está dentro de 2,3--3,1.

% ------------------------------------------------------------------------------
\subsection{Balance de volúmenes absolutos}
\label{sec:con-balance}

\ruta{service/VolumenAbsolutoService.\allowbreak{}yefriLermaCalcular(Datos)} completa un metro cúbico de concreto con agregado fino a
partir de parámetros elegidos por el profesional: agua $W$ (kg/m$^3$), relación agua/cemento $a/c$, densidades
relativas del cemento $\rho_c$, del grueso SSS $\rho_g$ y del fino SSS $\rho_f$, masa de grueso SSS $G$ (kg/m$^3$) y aire
$\alpha$ (\%). Con la densidad del agua tomada como 1000\,kg/m$^3$:

\begin{align}
	C &= \frac{W}{a/c}, \\
	V_w &= \frac{W}{1000}, \qquad V_c = \frac{C}{1000\,\rho_c}, \qquad V_g = \frac{G}{1000\,\rho_g}, \qquad
	V_{\text{aire}} = \frac{\alpha}{100}, \\
	V_f &= 1 - V_w - V_c - V_g - V_{\text{aire}}, \\
	F &= 1000\,\rho_f\,V_f .
\end{align}

El servicio exige $W$, $a/c$, las tres densidades y $G$ positivos y finitos, $0 \le \alpha < 100$ y $V_f > 0$; si no
queda volumen para el fino, lanza «No queda volumen positivo para agregado fino».

\begin{table}[H]
	\centering
	\caption{Ejemplo de \ruta{VolumenAbsolutoTest.\allowbreak{}balanceaUnMetroCubico}.}
	\label{tab:con-ejemplo-balance}
	\begin{tabular}{L{4.4cm}L{4.4cm}c}
		\hline
		\textbf{Dato} & \textbf{Resultado} & \textbf{Volumen (m$^3$)} \\ \hline
		$W = 180$ kg, $a/c = 0{,}50$ & $C = 360$ kg & $V_w = 0{,}180$; $V_c = 0{,}120$ \\
		$\rho_c = 3{,}0$ & & \\
		$G = 1000$ kg, $\rho_g = 2{,}5$ & & $V_g = 0{,}400$ \\
		$\alpha = 2$\,\% & & $V_{\text{aire}} = 0{,}020$ \\
		$\rho_f = 2{,}5$ & $F = 700$ kg & $V_f = 0{,}280$ \\
		\hline
	\end{tabular}
\end{table}

\subsubsection{Pantalla y «Usar en una mezcla»}

\ruta{BalanceController} lee los siete campos (con coma decimal), muestra cemento, fino, agua, grueso y los cinco
volúmenes, y habilita «Usar en una mezcla». Cualquier cambio posterior en un campo invalida el resultado y vuelve a
deshabilitar el botón. Al pulsarlo, \ruta{onUsarEnMezcla} navega a «Mezclas y probetas»
(\ruta{MainController.\allowbreak{}fernandoHuancaIrA}) y llama a \ruta{MezclasController.\allowbreak{}yefriLermaNuevaMezclaDesde}, que abre
«Registrar dosificación» con:

\begin{itemize}
	\item nombre «Balance a/c 0.50» (la relación con dos decimales) y referencia «Balance de volúmenes absolutos»;
	\item cemento, agua de diseño, fino SSS y grueso SSS por m$^3$;
	\item una observación con la relación $a/c$ y el aire, que pide completar la $f'_c$, la humedad y la absorción.
\end{itemize}

La resistencia, la humedad y la absorción quedan vacías. Si no hay un proyecto activo, se avisa y el formulario no se
abre. Las dependencias entre los dos controladores se inyectan con \ruta{@Lazy} para evitar un ciclo.

% ------------------------------------------------------------------------------
\subsection{Diseño de mezcla ACI 211.1}
\label{sec:con-aci}

\ruta{service/DisenoMezclaACI.\allowbreak{}yefriLermaDisenar(Entrada)} calcula una dosificación de partida para concreto de peso
normal por el método de volumen absoluto de ACI~211.1, con la resistencia promedio requerida y los límites de
durabilidad de la E.060. Las masas resultantes de agregados están en condición saturada superficialmente seca (SSS).

\subsubsection{Datos de entrada}

\begin{table}[H]
	\centering
	\caption{Campos de \ruta{DisenoMezclaACI.Entrada} y sus validaciones.}
	\label{tab:con-entrada-aci}
	\begin{tabular}{L{5.4cm}L{8.2cm}}
		\hline
		\textbf{Dato} & \textbf{Regla} \\ \hline
		$f'_c$ (MPa), asentamiento (mm), TMA (mm), módulo de finura, peso unitario compactado del grueso, densidades &
		Finitos y positivos. \\
		Desviación estándar $s$ y n.º de ensayos $n$ & Opcionales ($s$ nulo y $n = 0$ si se dejan vacíos). \\
		Absorción del grueso & De 0 a 20\,\%. \\
		Densidades relativas de fino y grueso & De 2,0 a 3,2. \\
		Densidad del cemento & De 2,5 a 3,3\,g/cm$^3$. \\
		Peso unitario compactado del grueso & De 1000 a 2000\,kg/m$^3$. \\
		TMA & Al menos 9,5\,mm. \\
		Exposición, aire incorporado y severidad & Exposición \ruta{NORMAL} si es nula; severidad \ruta{MODERADA} si es
		nula. \\
		\hline
	\end{tabular}
\end{table}

\subsubsection{Paso 1: resistencia promedio requerida}

\ruta{yefriLermaFcr(fc, s, n, criterio)} aplica dos casos. \textbf{Sin registro estadístico} ($s$ nulo o $n < 15$):

\begin{equation}
	f'_{cr} =
	\begin{cases}
		f'_c + 7.0 & f'_c < 21\ \text{MPa},\\
		f'_c + 8.5 & 21 \le f'_c \le 35\ \text{MPa},\\
		f'_c + 9.8 & f'_c > 35\ \text{MPa}.
	\end{cases}
	\label{eq:con-fcr-sin}
\end{equation}

\textbf{Con desviación estándar} ($n \ge 15$, $s > 0$), con $s' = k(n)\,s$:

\begin{equation}
	f'_{cr} = \max\left\{ f'_c + 1.34\, s',\;
	\begin{array}{ll}
		f'_c + 2.33\, s' - 3.5 & \text{si } f'_c \le 35\ \text{MPa}\\
		0.9\, f'_c + 2.33\, s' & \text{si } f'_c > 35\ \text{MPa}
	\end{array}
	\right\} .
	\label{eq:con-fcr-con}
\end{equation}

\begin{table}[H]
	\centering
	\caption{Factor de corrección de la desviación estándar $k(n)$ (\ruta{yefriLermaFactorMuestra}).}
	\label{tab:con-factor}
	\begin{tabular}{cc}
		\hline
		\textbf{Número de ensayos $n$} & \textbf{$k(n)$} \\ \hline
		15 & 1,16 \\
		20 & 1,08 \\
		25 & 1,03 \\
		30 o más & 1,00 \\
		\hline
	\end{tabular}
\end{table}

Entre los valores de la tabla se interpola linealmente. El texto del criterio aplicado (\ruta{criterioFcr}) se
muestra con el resultado.

\begin{advertencia}[Verifíquese contra la norma]
	Para $f'_c > 35$\,MPa sin registro estadístico, el código suma 9,8\,MPa. En ACI~318, desde la edición de 2005, ese
	caso es $1{,}10\,f'_c + 5{,}0$\,MPa. Compruebe la tabla 5.3 de la E.060 vigente antes de usar este caso; la prueba
	\ruta{fcrSinRegistroEstadistico} fija el valor actual del código ($42 + 9{,}8$).
\end{advertencia}

\subsubsection{Paso 2: durabilidad}

La exposición elegida (\ruta{DisenoMezclaACI.Exposicion}) fija una relación $a/c$ máxima, una $f'_c$ mínima y si se
exige aire incorporado. Los valores son los del código; el propio servicio advierte que se tomaron de ACI~318 cap.~19 /
E.060 cap.~4 y que deben verificarse contra la edición vigente.

\begin{table}[H]
	\centering
	\caption{Exposiciones de \ruta{DisenoMezclaACI.Exposicion}.}
	\label{tab:con-exposicion}
	\begin{tabular}{L{6.6cm}ccc}
		\hline
		\textbf{Exposición} & \textbf{$a/c$ máx.} & \textbf{$f'_c$ mín. (MPa)} & \textbf{Aire} \\ \hline
		Normal (sin requisitos especiales) & --- & --- & no \\
		Servicio en agua: baja permeabilidad & 0,50 & 28 & no \\
		Congelamiento y deshielo en condición húmeda & 0,45 & 31 & sí \\
		Sulfatos moderados (SO$_4$ 0,10--0,20\,\%): cemento tipo II & 0,50 & 28 & no \\
		Sulfatos severos (SO$_4$ 0,20--2,00\,\%): cemento tipo V & 0,45 & 31 & no \\
		Sulfatos muy severos (SO$_4 > 2{,}00$\,\%): tipo V + puzolana & 0,45 & 31 & no \\
		Agua salobre o de mar / cloruros & 0,40 & 35 & no \\
		\hline
	\end{tabular}
\end{table}

Si la exposición exige aire, el diseño se calcula con aire incorporado aunque no se haya marcado, y se avisa. Si
$f'_c$ es menor que la mínima de la exposición, se avisa (no se cambia $f'_c$). También se avisa si $f'_c < 17$\,MPa,
mínimo que la E.060 exige para concreto estructural según el mensaje del código.

\subsubsection{Paso 3: relación agua/cemento}

La relación por resistencia $(a/c)_R$ se interpola en la tabla 6.3.4(a) de ACI~211.1 con $f'_{cr}$:

\begin{table}[H]
	\centering
	\caption{Relación $a/c$ en función de $f'_{cr}$ a 28 días (\ruta{YEFRI\_LERMA\_AC\_SIN\_AIRE} y \ruta{YEFRI\_LERMA\_AC\_CON\_AIRE}).}
	\label{tab:con-ac}
	\begin{tabular}{L{3.8cm}ccccccc}
		\hline
		\textbf{$f'_{cr}$ (MPa)} & 15 & 20 & 25 & 30 & 35 & 40 & 45 \\ \hline
		Sin aire incorporado & 0,79 & 0,69 & 0,61 & 0,54 & 0,47 & 0,42 & 0,38 \\
		Con aire incorporado & 0,70 & 0,60 & 0,52 & 0,45 & 0,39 & --- & --- \\
		\hline
	\end{tabular}
\end{table}

\begin{itemize}
	\item Sin aire, si $f'_{cr} > 45$\,MPa el diseño se rechaza (se remite a ACI~211.4 y a mezclas de prueba); con aire,
	      el límite es $f'_{cr} = 35$\,MPa.
	\item Si $f'_{cr} < 15$\,MPa, se usa la relación de 15\,MPa y se avisa.
\end{itemize}

La relación adoptada es la menor entre la de resistencia y la de durabilidad:

\begin{equation}
	a/c = \min\left\{ (a/c)_R,\; (a/c)_D \right\}, \qquad (a/c)_D = +\infty \text{ si la exposición no la limita}.
\end{equation}

Si rige la durabilidad, se avisa con ambos valores.

\subsubsection{Paso 4: tamaño máximo nominal, agua y aire}

\ruta{yefriLermaIndiceTma} toma el mayor TMA de la tabla que no supera el ingresado
($9{,}5;\ 12{,}5;\ 19;\ 25;\ 37{,}5;\ 50;\ 75;\ 150$\,mm) y avisa si difiere en más de 0,6\,mm. El agua de mezclado se
lee en la tabla 6.3.3 de ACI~211.1. El código asocia cada rango de asentamiento a su valor central (25--50\,mm
$\rightarrow$ 37,5; 75--100\,mm $\rightarrow$ 87,5; 150--175\,mm $\rightarrow$ 162,5) e \textbf{interpola linealmente}
entre esos centros. Fuera de 25--175\,mm usa el límite más cercano y avisa.

\begin{table}[H]
	\centering
	\caption{Agua de mezclado (kg/m$^3$) según asentamiento y TMA.}
	\label{tab:con-agua}
	\begin{tabular}{L{3.6cm}cccccccc}
		\hline
		\textbf{Asentamiento} & \textbf{9,5} & \textbf{12,5} & \textbf{19} & \textbf{25} & \textbf{37,5} & \textbf{50} &
		\textbf{75} & \textbf{150} \\ \hline
		\multicolumn{9}{l}{\textit{Sin aire incorporado}} \\
		25--50\,mm & 207 & 199 & 190 & 179 & 166 & 154 & 130 & 113 \\
		75--100\,mm & 228 & 216 & 205 & 193 & 181 & 169 & 145 & 124 \\
		150--175\,mm & 243 & 228 & 216 & 202 & 190 & 178 & 160 & --- \\
		\multicolumn{9}{l}{\textit{Con aire incorporado}} \\
		25--50\,mm & 181 & 175 & 168 & 160 & 150 & 142 & 122 & 107 \\
		75--100\,mm & 202 & 193 & 184 & 175 & 165 & 157 & 133 & 119 \\
		150--175\,mm & 216 & 205 & 197 & 184 & 174 & 166 & 154 & --- \\
		\hline
	\end{tabular}
\end{table}

Para TMA de 150\,mm la tabla no tiene el rango de 150--175\,mm; si el asentamiento supera 87,5\,mm se usa el valor de
75--100\,mm y se avisa.

\begin{table}[H]
	\centering
	\caption{Contenido de aire (\%) según TMA (\ruta{YEFRI\_LERMA\_AIRE\_ATRAPADO} y \ruta{YEFRI\_LERMA\_AIRE\_INCORPORADO}).}
	\label{tab:con-aire}
	\begin{tabular}{L{3.6cm}cccccccc}
		\hline
		\textbf{TMA (mm)} & \textbf{9,5} & \textbf{12,5} & \textbf{19} & \textbf{25} & \textbf{37,5} & \textbf{50} &
		\textbf{75} & \textbf{150} \\ \hline
		Atrapado (sin aire incorporado) & 3,0 & 2,5 & 2,0 & 1,5 & 1,0 & 0,5 & 0,3 & 0,2 \\
		Incorporado, exposición suave & 4,5 & 4,0 & 3,5 & 3,0 & 2,5 & 2,0 & 1,5 & 1,0 \\
		Incorporado, moderada & 6,0 & 5,5 & 5,0 & 4,5 & 4,5 & 4,0 & 3,5 & 3,0 \\
		Incorporado, severa & 7,5 & 7,0 & 6,0 & 6,0 & 5,5 & 5,0 & 4,5 & 4,0 \\
		\hline
	\end{tabular}
\end{table}

\subsubsection{Paso 5: cemento}

\begin{equation}
	C = \frac{W}{a/c}, \qquad \text{bolsas} = \frac{C}{42.5} .
\end{equation}

Se avisa si $C < 250$\,kg/m$^3$ (cohesión) o $C > 550$\,kg/m$^3$ (retracción y calor de hidratación).

\subsubsection{Paso 6: agregado grueso}

El volumen de grueso seco y compactado por unidad de volumen de concreto, $b$, se interpola en la tabla 6.3.6 de
ACI~211.1 con el módulo de finura de la arena, acotado a 2,40--3,00 (fuera de ese rango se usa el límite y se avisa):

\begin{table}[H]
	\centering
	\caption{Volumen de agregado grueso $b$ (\ruta{YEFRI\_LERMA\_VOL\_GRUESO}).}
	\label{tab:con-grueso}
	\begin{tabular}{ccccc}
		\hline
		\textbf{TMA (mm)} & \textbf{MF 2,40} & \textbf{MF 2,60} & \textbf{MF 2,80} & \textbf{MF 3,00} \\ \hline
		9,5 & 0,50 & 0,48 & 0,46 & 0,44 \\
		12,5 & 0,59 & 0,57 & 0,55 & 0,53 \\
		19 & 0,66 & 0,64 & 0,62 & 0,60 \\
		25 & 0,71 & 0,69 & 0,67 & 0,65 \\
		37,5 & 0,75 & 0,73 & 0,71 & 0,69 \\
		50 & 0,78 & 0,76 & 0,74 & 0,72 \\
		75 & 0,82 & 0,80 & 0,78 & 0,76 \\
		150 & 0,87 & 0,85 & 0,83 & 0,81 \\
		\hline
	\end{tabular}
\end{table}

\begin{equation}
	G_{\text{seco}} = b \cdot \mathrm{PUC}, \qquad G_{\text{SSS}} = G_{\text{seco}} \left(1 + \frac{\text{abs}_g}{100}\right),
\end{equation}

\noindent donde PUC es el peso unitario compactado del grueso (kg/m$^3$) y $\text{abs}_g$ su absorción (\%).

\subsubsection{Paso 7: agregado fino por volumen absoluto}

El fino se obtiene con el balance de la sección~\ref{sec:con-balance}, llamando a
\ruta{VolumenAbsolutoService.\allowbreak{}yefriLermaCalcular} con $W$, $a/c$, $\rho_c$, $G_{\text{SSS}}$, $\rho_g$, $\rho_f$ y el
aire del paso 4. Como el grueso entra en condición SSS, se usa su densidad relativa SSS.

\subsubsection{Interpolación}

Todas las tablas se leen con \ruta{yefriLermaInterpolar}, interpolación lineal por tramos con valores extremos
constantes fuera del rango:

\begin{equation}
	y(x) =
	\begin{cases}
		y_0 & x \le x_0,\\
		y_{i-1} + \dfrac{x - x_{i-1}}{x_i - x_{i-1}} \left(y_i - y_{i-1}\right) & x_{i-1} < x \le x_i,\\
		y_{n} & x > x_n.
	\end{cases}
\end{equation}

\subsubsection{Resultado y proporción en peso}

\ruta{Resultado} devuelve $f'_{cr}$ con su criterio, las relaciones $a/c$ de resistencia, de durabilidad y adoptada, el
TMA de tabla, el agua, el aire, el cemento y las bolsas, $b$, el grueso seco y SSS, el fino SSS, los cinco volúmenes y
la lista de avisos, que siempre termina recordando que es un cálculo de partida que debe verificarse con mezclas de
prueba. \ruta{proporcionEnPeso()} da la proporción respecto del cemento y el agua por bolsa:

\begin{equation}
	1 : \frac{F_{\text{SSS}}}{C} : \frac{G_{\text{SSS}}}{C} \;\big/\; 42.5 \cdot (a/c)\ \text{L de agua por bolsa de 42,5 kg}.
\end{equation}

\subsubsection{Ejemplo de cálculo}

Con los datos de la prueba \ruta{DisenoMezclaTest.\allowbreak{}ejemploFc28SinAire}: $f'_c = 28$\,MPa sin registro estadístico,
asentamiento 87,5\,mm, TMA 19\,mm, MF 2,7, PUC 1600\,kg/m$^3$, $\rho_c = 3{,}15$, $\rho_f = \rho_g = 2{,}65$, absorción
del grueso 1\,\%, exposición normal y sin aire incorporado.

\begin{table}[H]
	\centering
	\caption{Diseño de la prueba \ruta{ejemploFc28SinAire} (las cifras marcadas con * se calcularon a partir del código).}
	\label{tab:con-ejemplo-aci}
	\begin{tabular}{L{4.6cm}L{5.4cm}L{3.2cm}}
		\hline
		\textbf{Paso} & \textbf{Cálculo} & \textbf{Resultado} \\ \hline
		$f'_{cr}$ & $28 + 8{,}5$ & 36,5\,MPa \\
		$a/c$ por resistencia & $0{,}47 - 0{,}05 \cdot (1{,}5/5)$ & 0,455 \\
		Agua & tabla, 75--100\,mm y 19\,mm & 205\,kg/m$^3$ \\
		Aire atrapado & tabla, 19\,mm & 2,0\,\% \\
		Cemento & $205 / 0{,}455$ & 450,5\,kg/m$^3$ (10,6 bolsas)* \\
		$b$ & entre MF 2,6 (0,64) y 2,8 (0,62) & 0,63 \\
		Grueso seco y SSS & $0{,}63 \cdot 1600$; $\times 1{,}01$ & 1008 y 1018,1\,kg/m$^3$* \\
		Volúmenes & agua, cemento, grueso, aire & 0,2050; 0,1430; 0,3842; 0,0200* \\
		Fino & $V_f = 0{,}2478$; $\times 2650$ & 656,6\,kg/m$^3$* \\
		Proporción & & 1 : 1,46 : 2,26 / 19,3\,L* \\
		\hline
	\end{tabular}
\end{table}

La prueba verifica $f'_{cr}$, $(a/c)_R$, el agua, el aire, el cemento, $b$, el grueso seco, que la suma de los cinco
volúmenes sea 1 y que el fino esté entre 500 y 900\,kg/m$^3$.

\subsubsection{Cemento frente a sulfatos}

\ruta{yefriLermaAvisoCemento(exposicion, tipoMaterial)} compara el tipo del cemento elegido en la biblioteca con la
exposición. Primero quita del tipo el prefijo «CEMENTO» y, si está, «TIPO» (así «CEMENTO TIPO IP» queda en «IP» y
«CEMENTO» queda vacío).

\begin{table}[H]
	\centering
	\caption{Verificación del tipo de cemento.}
	\label{tab:con-sulfatos}
	\begin{tabular}{L{3.4cm}L{3.8cm}L{6.4cm}}
		\hline
		\textbf{Exposición} & \textbf{Tipos aceptados} & \textbf{Aviso} \\ \hline
		Sulfatos moderados & II, MS, V, HS & «se requiere tipo II o MS» si no es aceptado. \\
		Sulfatos severos & V, HS & «se requiere tipo V o HS». \\
		Sulfatos muy severos & V, HS & Siempre avisa: con V o HS, que además se requiere puzolana; si no, que se requiere
		tipo V o HS más puzolana. \\
		Otras & cualquiera & Sin aviso. \\
		\hline
	\end{tabular}
\end{table}

Un cemento genérico se informa como «sin tipo indicado». Si no se eligió un cemento y la exposición es de sulfatos, la
pantalla pide elegirlo para poder verificarlo.

\subsubsection{Flujo del asistente}

\ruta{MezclasController}:

\begin{enumerate}
	\item «Diseñar mezcla» (\ruta{onDisenar}) abre el panel con los combos de agregado fino, agregado grueso y cemento
	      llenos desde la biblioteca. La exposición empieza en «Normal» y la severidad en «Moderada» (habilitada solo
	      si se marca «Aire incorporado»).
	\item Elegir un agregado fino copia su módulo de finura y su densidad relativa; uno grueso, su tamaño máximo, su
	      peso unitario compactado, su densidad relativa y su absorción; un cemento, su densidad relativa como densidad
	      en g/cm$^3$. Los campos pueden escribirse a mano.
	\item «Calcular diseño» (\ruta{onCalcularDiseno}) marca los campos obligatorios vacíos (todos salvo $s$ y $n$),
	      redondea $n$ al entero más cercano, llama a \ruta{yefriLermaDisenar}, agrega el aviso de cemento y habilita
	      «Usar en una mezcla».
	\item «Usar en una mezcla» (\ruta{onUsarDiseno}) exige un proyecto activo y abre «Registrar dosificación» con el
	      nombre «Diseño ACI $f'_c$ \textit{valor} MPa» (la resistencia con tres decimales), la referencia «Asistente
	      ACI 211.1 + E.060», la $f'_c$, el cemento, el agua, el fino SSS y el grueso SSS. Si se eligieron agregados de la biblioteca, también los selecciona y copia
	      su absorción y humedad; sin grueso de biblioteca, copia la absorción escrita en el asistente. Las observaciones
	      reciben todos los avisos del diseño.
\end{enumerate}

% ------------------------------------------------------------------------------
\subsection{Mezclas, tanda y corrección por humedad}
\label{sec:con-mezclas}

\subsubsection{Registro de la dosificación}

Una \ruta{Mezcla} guarda, por m$^3$, el cemento $C$, el agua de diseño $W$ y las masas SSS de fino $F$ y grueso $G$,
además de la $f'_c$ especificada, la absorción y la humedad total de cada agregado, la referencia, las observaciones y
el texto de los materiales elegidos («nombre \textbar\ informe: \dots\ \textbar\ fecha: \dots»).
\ruta{MezclaService.\allowbreak{}yefriLermaGuardar} exige un proyecto existente, nombre, $f'_c$, $C$, $W$, $F$ y $G$ positivos y
finitos, absorciones y humedades entre 0 y 100\,\%, y además que la corrección para 1\,m$^3$ sea posible (no puede
quedar agua a agregar negativa). Las mezclas guardadas con un esquema anterior, sin las columnas de materiales, se leen
con esos campos nulos (prueba \ruta{CompatibilidadMezclaTest}).

\subsubsection{Corrección por humedad}

\ruta{MezclaService.\allowbreak{}yefriLermaCalcular(Mezcla, V)} convierte la dosificación SSS a las cantidades de una tanda de
volumen $V$ con los agregados en su humedad real. La humedad $w$ y la absorción $\text{abs}$ son porcentajes de la masa
seca:

\begin{align}
	S_f &= \frac{F}{1 + \text{abs}_f/100}, \qquad S_g = \frac{G}{1 + \text{abs}_g/100} && \text{(masas secas)}, \\
	F_h &= S_f \left(1 + \frac{w_f}{100}\right), \qquad G_h = S_g \left(1 + \frac{w_g}{100}\right) && \text{(masas húmedas)}, \\
	W_{\text{agregar}} &= W - S_f\,\frac{w_f - \text{abs}_f}{100} - S_g\,\frac{w_g - \text{abs}_g}{100}
	&& \text{(agua libre)}.
\end{align}

Las cantidades de la tanda son $C\,V$, $F_h\,V$, $G_h\,V$ y $W_{\text{agregar}}\,V$. La relación $a/c$ que se informa es
la de diseño, $W/C$. El cálculo rechaza $V \le 0$, un agua a agregar negativa («El agua libre supera el agua de diseño»)
y resultados no finitos.

\begin{table}[H]
	\centering
	\caption{Ejemplo de \ruta{MezclaIntegracionTest.\allowbreak{}registraYCalculaConHumedadYEnsayo} ($V = 2$\,m$^3$).}
	\label{tab:con-ejemplo-humedad}
	\begin{tabular}{L{3cm}L{3.4cm}L{3.8cm}L{3cm}}
		\hline
		\textbf{Material} & \textbf{Diseño (por m$^3$)} & \textbf{Seco y húmedo (por m$^3$)} & \textbf{Tanda de 2\,m$^3$} \\ \hline
		Cemento & 350\,kg & --- & 700\,kg \\
		Fino ($\text{abs}=1$, $w=3$) & 808\,kg SSS & 800 y 824\,kg & 1648\,kg \\
		Grueso ($\text{abs}=1$, $w=0{,}5$) & 1010\,kg SSS & 1000 y 1005\,kg & 2010\,kg \\
		Agua & 180\,kg & $180 - 16 + 5 = 169$\,kg & 338\,kg \\
		\hline
	\end{tabular}
\end{table}

\subsubsection{Volumen de la tanda}

El volumen se escribe en «Volumen del lote (m$^3$)» (1 por defecto). «Elegir elementos del modelo»
(\ruta{onElegirElementos} y \ruta{onUsarSeleccion}) permite seleccionar elementos del proyecto y escribe la suma de sus
volúmenes \textbf{brutos} ($l\,a\,h$); la propia pantalla advierte que esa suma no descuenta intersecciones y remite al
metrado. Al seleccionar una mezcla, el panel de resultados muestra la tanda y, debajo, la evaluación E.060 de sus
probetas (sección~\ref{sec:con-ensayos}).

% ------------------------------------------------------------------------------
\subsection{Probetas y evaluación de resistencia}
\label{sec:con-ensayos}

\subsubsection{Registro de probetas}

\ruta{MezclaService.\allowbreak{}yefriLermaGuardarEnsayo} guarda un \ruta{Ensayo} si la mezcla existe, la identificación no está
vacía, la fecha es una fecha ISO válida, la edad y la resistencia son positivas y el asentamiento no es negativo.

\subsubsection{Criterios de aceptación}

\ruta{service/EvaluacionResistencia.\allowbreak{}yefriLermaEvaluar(fc, resultados)} aplica los dos criterios de aceptación que el
código atribuye a la E.060 (art.~5.6) y ACI~318. Con los resultados $x_1, \dots, x_n$ (se descartan los nulos, no
finitos o no positivos):

\begin{align}
	\text{(a)} \quad & \frac{x_i + x_{i+1} + x_{i+2}}{3} \ge f'_c \qquad \text{para todo } i = 1, \dots, n-2, \\
	\text{(b)} \quad & \min_i x_i \ge x_{\lim}, \qquad
	x_{\lim} =
	\begin{cases}
		f'_c - 3.5 & f'_c \le 35\ \text{MPa},\\
		0.9\, f'_c & f'_c > 35\ \text{MPa}.
	\end{cases}
\end{align}

El resultado es \ruta{conforme()} solo si hay al menos tres resultados (\ruta{completo}) y se cumplen (a) y (b). Con
menos de tres, el mensaje indica que falta el criterio de promedios y, si algún resultado ya está bajo $x_{\lim}$, lo
advierte. El servicio calcula también el promedio, el mínimo, el menor promedio de tres consecutivos y, con $n \ge 2$, la
desviación estándar muestral

\begin{equation}
	s = \sqrt{\frac{1}{n-1} \sum_{i=1}^{n} \left(x_i - \bar{x}\right)^2 } .
\end{equation}

\subsubsection{Qué resultados se evalúan}

\ruta{MezclasController.\allowbreak{}yefriLermaEvaluacion} toma solo las probetas con edad de $28 \pm 0{,}5$ días y las ordena de la
más antigua a la más reciente según el orden de registro (la lista del servicio viene del más reciente al más antiguo y
el controlador la recorre al revés). Cada probeta cuenta como un resultado.

\begin{importante}[Probetas y ensayos]
	En la norma, un «ensayo» es el promedio de al menos dos probetas de la misma muestra a la edad de diseño. La
	aplicación evalúa cada probeta registrada como un ensayo, por lo que el resultado es orientativo si se registran
	probetas sueltas; así lo advierte el comentario de la clase.
\end{importante}

\begin{table}[H]
	\centering
	\caption{Casos de \ruta{DisenoMezclaTest.\allowbreak{}evaluacionDeResistencia}.}
	\label{tab:con-ejemplo-evaluacion}
	\begin{tabular}{cL{4.2cm}L{7.4cm}}
		\hline
		\textbf{$f'_c$} & \textbf{Resultados (MPa)} & \textbf{Evaluación} \\ \hline
		21 & 23; 22; 21,5; 24 & Conforme. \\
		21 & 21; 20,5; 20; 25 & No cumple (a): el primer promedio es 20,5. \\
		21 & 30; 30; 17 & No cumple (b): $17 < 21 - 3{,}5 = 17{,}5$. \\
		21 & 25; 24 & Incompleto: no conforme por tener menos de tres. \\
		40 & 41 & Límite individual $0{,}9 \cdot 40 = 36$. \\
		\hline
	\end{tabular}
\end{table}

% ------------------------------------------------------------------------------
\subsection{Reporte PDF}
\label{sec:con-pdf}

«Guardar PDF» llama a \ruta{ReporteMezcla.\allowbreak{}yefriLermaGenerar(proyecto, mezcla, V, ensayos)}, que genera con OpenPDF
(\ruta{com.lowagie.text}) un documento A4 con márgenes de 40\,pt y el siguiente contenido:

\begin{enumerate}
	\item título «Resumen de dosificación», proyecto, mezcla, referencia y materiales fino y grueso;
	\item datos de diseño ingresados: $f'_c$, cemento, agua de diseño, fino y grueso SSS con su absorción y humedad, y
	      observaciones;
	\item cantidades corregidas para el volumen $V$ (\ruta{MezclaService.\allowbreak{}yefriLermaCalcular}): cemento, agua a añadir,
	      fino y grueso húmedos y relación $a/c$ de diseño;
	\item lista de ensayos físicos (identificación, fecha, edad, resistencia, asentamiento y observaciones);
	\item una nota final que aclara que es un registro de la dosificación proporcionada por el usuario, que no
	      constituye diseño ACI, aprobación ni verificación de conformidad, y que requiere revisión profesional y mezcla
	      de prueba.
\end{enumerate}

Los números se escriben con tres decimales y punto decimal (\ruta{Locale.ROOT}). El PDF no incluye la evaluación E.060
ni el detalle del asistente ACI (salvo lo que se haya copiado en la referencia y las observaciones).

% ------------------------------------------------------------------------------
\subsection{Pruebas automatizadas}
\label{sec:con-pruebas}

\begin{longtable}{L{4.6cm}L{9cm}}
	\caption{Pruebas que cubren este capítulo (\ruta{src/test/java/com/concretolab}).} \label{tab:con-pruebas} \\
	\hline
	\textbf{Clase} & \textbf{Qué verifica} \\ \hline
	\endfirsthead
	\hline
	\textbf{Clase} & \textbf{Qué verifica} \\ \hline
	\endhead
	\hline
	\endfoot
	\ruta{DisenoMezclaTest} & $f'_{cr}$ sin registro y con desviación (incluido el factor para 20 ensayos y el caso
	$f'_c > 35$); ejemplo completo de 28\,MPa (tabla~\ref{tab:con-ejemplo-aci}); durabilidad que gobierna la relación
	$a/c$; congelamiento que fuerza aire incorporado; rechazo fuera de tabla; evaluación de resistencia
	(tabla~\ref{tab:con-ejemplo-evaluacion}); huso fino y módulo de finura (tabla~\ref{tab:con-ejemplo-granulometria}). \\
	\ruta{PendientesCierreTest} & Catálogo de tipos de cemento NTP y aditivo sin perder los tipos anteriores; aviso de
	cemento frente a sulfatos (tabla~\ref{tab:con-sulfatos}); además, las cuantías de acero del metrado. \\
	\ruta{GranulometriaTest} & Porcentajes, acumulados, \% que pasa y balance de masas; rechazo de aberturas desordenadas
	y masas imposibles. \\
	\ruta{VolumenAbsolutoTest} & Balance de un metro cúbico (tabla~\ref{tab:con-ejemplo-balance}) y rechazo de un
	volumen inviable. \\
	\ruta{MezclaIntegracionTest} & Registro con base SQLite temporal, corrección por humedad
	(tabla~\ref{tab:con-ejemplo-humedad}), ensayo y texto del PDF; materiales y granulometrías guardados y datos
	inválidos rechazados. \\
	\ruta{CompatibilidadMezclaTest} & Lectura de una mezcla creada con el esquema anterior de la tabla \ruta{mezclas}. \\
\end{longtable}

% ------------------------------------------------------------------------------
\subsection{Supuestos y límites}
\label{sec:con-limites}

\begin{itemize}
	\item \textbf{Alcance del diseño}: concreto de peso normal con agregado grueso angular y sin aditivos reductores de
	      agua (las tablas de ACI~211.1 lo suponen). No cubre concreto liviano ni de alta resistencia: se rechaza
	      $f'_{cr} > 45$\,MPa (35\,MPa con aire incorporado). Los aditivos registrados en la biblioteca no intervienen en
	      ningún cálculo.
	\item \textbf{Tablas e interpolación}: el agua se interpola entre los centros de los rangos de asentamiento, una
	      elección del código que la tabla original no prescribe; el TMA se redondea al inmediato inferior de la tabla y
	      el módulo de finura se acota a 2,40--3,00.
	\item \textbf{$f'_{cr}$ sin registro estadístico para $f'_c > 35$\,MPa}: el código usa $f'_c + 9{,}8$; verifíquese
	      contra la tabla 5.3 de la E.060 (véase el recuadro del paso 1).
	\item \textbf{Durabilidad}: las relaciones $a/c$ máximas, las $f'_c$ mínimas y los rangos de sulfatos de la
	      tabla~\ref{tab:con-exposicion} son los del código; el propio servicio pide verificarlos contra la edición
	      vigente de la E.060 y de ACI~318. Solo se verifica el tipo de cemento frente a sulfatos, y solo si se elige un
	      cemento de la biblioteca.
	\item \textbf{Densidades}: la densidad del agua se toma como 1000\,kg/m$^3$ y la densidad del cemento se usa como
	      densidad relativa. En el diseño, la absorción del fino no interviene (el fino se obtiene directamente en
	      condición SSS).
	\item \textbf{Humedad}: la corrección supone que la humedad total y la absorción son porcentajes de la masa seca y
	      que las masas de diseño están en condición SSS. La proporción «L de agua por bolsa» usa la relación $a/c$ de
	      diseño, no el agua corregida.
	\item \textbf{Granulometría}: los porcentajes se calculan sobre la masa inicial; el módulo de finura exige seis
	      tamices cualesquiera de la serie estándar, aunque el comentario del código dice que para agregado fino se exigen
	      los de 0,15 a 4,75\,mm. Los límites de los husos son los del código; verifíquense contra la edición de la
	      NTP~400.037 y de la ASTM~C33 que corresponda. El módulo de finura calculado no actualiza el material.
	\item \textbf{Evaluación}: cada probeta de $28 \pm 0{,}5$ días cuenta como un ensayo; el orden es el de registro, no
	      el de la fecha; la desviación estándar se calcula pero no se muestra en el mensaje.
	\item \textbf{Volumen desde el modelo}: la suma de elementos usa volúmenes brutos; para un volumen sin cruces use el
	      metrado neto del modelo.
	\item \textbf{Reporte}: el PDF siempre declara que no constituye un diseño ACI, incluso si la dosificación provino
	      del asistente, y no incluye la evaluación E.060.
	\item \textbf{Biblioteca}: los materiales no se editan ni se eliminan desde la pantalla; las correcciones se
	      registran como materiales nuevos.
\end{itemize}

\begin{nota}[Responsabilidad profesional]
	Todos los resultados de este módulo son valores de partida o de control. La dosificación definitiva debe
	validarse con mezclas de prueba y ensayos de laboratorio, y la conformidad la dictamina el profesional responsable.
\end{nota}
